\documentclass{article}
\usepackage{amsmath,amssymb,amsthm}
\usepackage{algorithm}
\usepackage{algorithmic}
\usepackage{bm}
\usepackage{hyperref}
\newtheorem{theorem}{Theorem}
\newtheorem{lemma}{Lemma}
\newtheorem{assumption}{Assumption}
\begin{document}

\section*{Clipped‑Momentum Gradient Descent (CMGD)}

\section*{Mathematical Formulation}
Let $f:\mathbb{R}^d\rightarrow\mathbb{R}$ be a (possibly non‑convex) differentiable objective.  
At iteration $t\ge 0$ we maintain a parameter vector $x_t\in\mathbb{R}^d$ and a momentum buffer $m_t\in\mathbb{R}^d$.
Given hyper‑parameters
\[
\eta>0\;\text{(stepsize)},\qquad
\beta\in[0,1)\;\text{(momentum coefficient)},\qquad
C>0\;\text{(gradient‑clipping threshold)},
\]
the updates are

\begin{align}
g_t^{\text{raw}} &:= \nabla f(x_t) \quad\text{(true gradient or unbiased stochastic estimate)}\label{eq:raw}\\[4pt]
\tilde g_t &:= \frac{g_t^{\text{raw}}}{\displaystyle \max\!\Bigl(1,\;\frac{\|g_t^{\text{raw}}\|}{C}\Bigr)} 
            \;=\;
            \begin{cases}
                g_t^{\text{raw}}, & \|g_t^{\text{raw}}\|\le C,\\[4pt]
                C\,\dfrac{g_t^{\text{raw}}}{\|g_t^{\text{raw}}\|}, & \|g_t^{\text{raw}}\|> C,
            \end{cases}\label{eq:clip}\\[6pt]
m_t &:= \beta\,m_{t-1} + (1-\beta)\,\tilde g_t,\qquad m_{-1}=0,\label{eq:momentum}\\[4pt]
x_{t+1} &:= x_t - \eta\, m_t.\label{eq:update}
\end{align}

Thus every raw gradient is first clipped to have Euclidean norm at most $C$, then a standard exponential moving‑average momentum is applied, and finally a gradient step is taken.

\section*{Pseudocode}
\begin{algorithm}[H]
\caption{Clipped‑Momentum Gradient Descent (CMGD)}
\begin{algorithmic}[1]
\STATE {\bf Input:} stepsize $\eta>0$, momentum $\beta\in[0,1)$, clipping threshold $C>0$, 
maximum iterations $T$.
\STATE {\bf Initialize:} $x_0\in\mathbb{R}^d$, $m_{-1}=0$.
\FOR{$t=0,\dots,T-1$}
    \STATE Compute (or sample) raw gradient $g_t^{\text{raw}}=\nabla f(x_t)$.
    \STATE Clip:
    \[
    \tilde g_t \leftarrow 
    \frac{g_t^{\text{raw}}}{\max\!\bigl(1,\|g_t^{\text{raw}}\|/C\bigr)}.
    \]
    \STATE Update momentum:
    \[
    m_t \leftarrow \beta\,m_{t-1}+(1-\beta)\,\tilde g_t .
    \]
    \STATE Gradient step:
    \[
    x_{t+1}\leftarrow x_t-\eta\,m_t .
    \]
\ENDFOR
\STATE {\bf Output:} $x_T$ (or the averaged iterate $\bar x_T:=\frac1T\sum_{t=0}^{T-1}x_t$).
\end{algorithmic}
\end{algorithm}

\section*{Dry Run Example}
Consider the univariate quadratic $f(w)=(w-3)^2$, with exact gradients.
Set $w_0=0$, stepsize $\eta=0.1$, momentum $\beta=0.5$, clipping threshold $C=1$.
We run three iterations.

\begin{center}
\begin{tabular}{c|c|c|c|c}
$t$ & $g_t^{\text{raw}}$ & $\tilde g_t$ (clipped) & $m_t$ & $w_{t+1}$ \\
\hline
0 & $2(w_0-3) = -6$ & $C\frac{-6}{|{-6}|}=-1$ & $0.5\cdot0+0.5(-1)=-0.5$ & $0-0.1(-0.5)=0.05$ \\
1 & $2(0.05-3)=-5.9$ & $-1$ & $0.5(-0.5)+0.5(-1)=-0.75$ & $0.05-0.1(-0.75)=0.125$ \\
2 & $2(0.125-3)=-5.75$ & $-1$ & $0.5(-0.75)+0.5(-1)=-0.875$ & $0.125-0.1(-0.875)=0.2125$
\end{tabular}
\end{center}

Objective values:
\[
f(w_0)=9,\qquad f(w_1)= (0.05-3)^2\approx 8.7025,\qquad
f(w_2)= (0.125-3)^2\approx 8.2656,\qquad
f(w_3)= (0.2125-3)^2\approx 7.7490.
\]
The iterates move toward the minimiser $w^\star=3$, while each gradient is clipped to norm~$1$.

\newpage
\section*{Assumptions}
\begin{assumption}[Smoothness]\label{as:smooth}
$f$ is $L$‑smooth: for all $x,y\in\mathbb{R}^d$,
\[
\|\nabla f(x)-\nabla f(y)\|\le L\|x-y\|.
\tag{A1}
\]
Equivalently,
\[
f(y)\le f(x)+\langle\nabla f(x),y-x\rangle+\frac{L}{2}\|y-x\|^2 .
\tag{A1'}
\]
\end{assumption}

\begin{assumption}[Lower Boundedness]\label{as:lb}
There exists $f_{\inf}>-\infty$ such that $\forall x,\;f(x)\ge f_{\inf}$. Define $\Delta:=f(x_0)-f_{\inf}<\infty$.
\tag{A2}
\end{assumption}

\begin{assumption}[Stochastic Gradient Unbiasedness and Bounded Variance]\label{as:stoch}
We allow a stochastic estimator $g_t$ of the raw gradient satisfying
\[
\mathbb{E}[g_t \mid \mathcal{F}_t]=\nabla f(x_t),\qquad
\mathbb{E}[\|g_t-\nabla f(x_t)\|^2\mid \mathcal{F}_t]\le\sigma^2,
\tag{A3}
\]
where $\mathcal{F}_t:=\sigma\{x_0,\dots,x_t\}$.
\end{assumption}

\begin{assumption}[Clipping Threshold]\label{as:clip}
The clipping constant $C>0$ is fixed and finite.
\tag{A4}
\end{assumption}

\begin{assumption}[Momentum Parameter]\label{as:beta}
$\beta\in[0,1)$.
\tag{A5}
\end{assumption}

All subsequent analysis holds under \ref{as:smooth}–\ref{as:beta}.

\section*{Main Convergence Theorem}
\begin{theorem}[Non‑convex convergence of CMGD]\label{thm:main}
Let $\{x_t\}_{t\ge0}$ be generated by CMGD with stepsize $\eta\le \dfrac{1}{L(1+\beta)}$.
Define the averaged iterate $\bar x_T:=\frac{1}{T}\sum_{t=0}^{T-1}x_t$.
Then for any $T\ge1$,
\[
\frac{1}{T}\sum_{t=0}^{T-1}\mathbb{E}\bigl[\|\nabla f(x_t)\|^2\bigr]
\;\le\;
\underbrace{\frac{2\Delta}{\eta(1-\beta)T}}_{\text{optimization term}}
\;+\;
\underbrace{2L\eta C^2}_{\text{clipping‑induced bias}}.
\tag{1}
\]
Consequently, choosing $\eta = \Theta\!\bigl(\frac{1}{\sqrt{T}}\bigr)$ yields
\[
\frac{1}{T}\sum_{t=0}^{T-1}\mathbb{E}\bigl[\|\nabla f(x_t)\|^2\bigr]=
O\!\bigl(T^{-1/2}\bigr).
\]
\end{theorem}

\section*{Theoretical Properties}
\begin{itemize}
\item \textbf{Stability.} Clipping guarantees $\|\tilde g_t\|\le C$ a.s., which prevents exploding updates even when $\|\nabla f\|$ is large.
\item \textbf{Hyperparameter Sensitivity.}
    \begin{enumerate}
    \item Stepsize $\eta$ must respect $\eta\le \frac{1}{L(1+\beta)}$ to keep the descent inequality \eqref{eq:descent}.
    \item Larger $C$ reduces the bias term $2L\eta C^2$ but allows larger raw gradients; a typical choice is $C\approx \sqrt{\Delta/L}$.
    \item Momentum $\beta$ appears only through the factor $(1-\beta)$ in the optimisation term, i.e., higher momentum slows the decay of the first term but can improve practical speed.
    \end{enumerate}
\item \textbf{Comparison to Baselines.}
    \begin{itemize}
    \item Without clipping, standard SGD with momentum enjoys the same $O(T^{-1/2})$ rate but requires a bounded gradient assumption, which may be violated for non‑convex deep nets.
    \item Compared with Adam, CMGD is simpler; the convergence rate matches the best known $O(\log T/T)$ for Adam‑type methods under similar assumptions, while the proof is elementary because clipping supplies a uniform gradient bound.
    \end{itemize}
\end{itemize}

\section*{Discussion}
The theorem shows that gradient clipping converts an arbitrary smooth non‑convex problem into one where the effective gradients are uniformly bounded by $C$. This enables a classic descent‑lemma argument leading to the $O(T^{-1/2})$ rate. The extra term $2L\eta C^2$ reflects the bias introduced by clipping: when the true gradient norm exceeds $C$, the algorithm uses a scaled version, thus deviating from the true descent direction. By letting $\eta$ decay as $1/\sqrt{T}$ the bias term vanishes at the same rate as the optimisation term, yielding the optimal $O(T^{-1/2})$ convergence for stochastic first‑order methods on smooth non‑convex objectives.

\newpage
\section*{Preliminaries (Definitions and Lemmas)}
\begin{lemma}[Descent Lemma]\label{lem:descent}
If $f$ is $L$‑smooth, then for any $x$ and any vector $d$,
\[
f(x-d)\le f(x)-\langle\nabla f(x),d\rangle+\frac{L}{2}\|d\|^2 .
\tag{2}
\]
\end{lemma}
\begin{proof}
Directly from \ref{as:smooth} (inequality \eqref{eq:raw}) applied with $y=x-d$.
\end{proof}

\begin{lemma}[Momentum Bound]\label{lem:momentum}
For any sequence $\{\tilde g_t\}$ with $\|\tilde g_t\|\le C$, the momentum iterates satisfy
\[
\|m_t\|\le C,\qquad\forall t\ge0.
\tag{3}
\]
\end{lemma}
\begin{proof}
Induction. Base $m_{-1}=0$ satisfies $\|m_{-1}\|\le C$. Assume $\|m_{t-1}\|\le C$.
Then
\[
\|m_t\| =\|\beta m_{t-1}+(1-\beta)\tilde g_t\|
\le \beta\|m_{t-1}\|+(1-\beta)\|\tilde g_t\|
\le \beta C+(1-\beta)C=C.
\]
\end{proof}

\section*{Main Proof (Step‑by‑Step Derivation)}
We prove Theorem~\ref{thm:main}. Throughout, $\mathcal{F}_t$ denotes the filtration generated by $\{x_0,\dots,x_t\}$.

\noindent\textbf{[Step 1] Apply the descent lemma to the update \eqref{eq:update}.}
\[
f(x_{t+1})\le f(x_t)-\langle\nabla f(x_t),\eta m_t\rangle
      +\frac{L}{2}\eta^2\|m_t\|^2 .
\tag{4}
\]

\noindent\textbf{[Step 2] Replace $m_t$ using its definition \eqref{eq:momentum}.}
\[
\langle\nabla f(x_t),m_t\rangle
= \beta\langle\nabla f(x_t),m_{t-1}\rangle
  +(1-\beta)\langle\nabla f(x_t),\tilde g_t\rangle .
\tag{5}
\]

\noindent\textbf{[Step 3] Take conditional expectation $\mathbb{E}[\cdot\mid\mathcal{F}_t]$.
Use unbiasedness of the raw stochastic gradient (Assumption~\ref{as:stoch}) and the deterministic clipping operator.}
Because $\tilde g_t$ is a deterministic function of $g_t^{\text{raw}}$,
\[
\mathbb{E}\bigl[\tilde g_t\mid\mathcal{F}_t\bigr]
= \mathbb{E}\!\left[\frac{g_t^{\text{raw}}}{\max(1,\|g_t^{\text{raw}}\|/C)}\Bigm|\mathcal{F}_t\right].
\]
Define the \emph{clipping bias} 
\[
b_t:=\nabla f(x_t)-\mathbb{E}[\tilde g_t\mid\mathcal{F}_t].
\tag{6}
\]
Then
\[
\mathbb{E}[\langle\nabla f(x_t),\tilde g_t\rangle\mid\mathcal{F}_t]
= \|\nabla f(x_t)\|^2-\langle\nabla f(x_t),b_t\rangle .
\tag{7}
\]

\noindent\textbf{[Step 4] Bound the bias term.}
Since $\|\tilde g_t\|\le C$ (by construction) and $\mathbb{E}[\tilde g_t\mid\mathcal{F}_t]$ is a convex combination of vectors with norm $\le C$, we have
\[
\|\mathbb{E}[\tilde g_t\mid\mathcal{F}_t]\|\le C.
\]
Thus,
\[
\|b_t\|=\|\nabla f(x_t)-\mathbb{E}[\tilde g_t\mid\mathcal{F}_t]\|
\le \|\nabla f(x_t)\|+C .
\]
Using Cauchy–Schwarz,
\[
\langle\nabla f(x_t),b_t\rangle\le \|\nabla f(x_t)\|\,\|b_t\|
\le \|\nabla f(x_t)\|(\|\nabla f(x_t)\|+C)
= \|\nabla f(x_t)\|^2 + C\|\nabla f(x_t)\|.
\tag{8}
\]

\noindent\textbf{[Step 5] Plug (5)–(8) into (4) and take total expectation.}
First, conditionally on $\mathcal{F}_t$,
\[
\begin{aligned}
\mathbb{E}[f(x_{t+1})\mid\mathcal{F}_t]
&\le f(x_t)-\eta\Bigl[
\beta\langle\nabla f(x_t),m_{t-1}\rangle
+(1-\beta)\bigl(\|\nabla f(x_t)\|^2-\langle\nabla f(x_t),b_t\rangle\bigr)
\Bigr]
+\frac{L}{2}\eta^2\|m_t\|^2 .
\end{aligned}
\tag{9}
\]

\noindent\textbf{[Step 6] Upper‑bound the term involving $m_{t-1}$.}
Using Lemma~\ref{lem:momentum} ($\|m_{t-1}\|\le C$) and Cauchy–Schwarz,
\[
\langle\nabla f(x_t),m_{t-1}\rangle\le \|\nabla f(x_t)\|\,C .
\tag{10}
\]

\noindent\textbf{[Step 7] Combine (9), (8) and (10).}
\[
\begin{aligned}
\mathbb{E}[f(x_{t+1})\mid\mathcal{F}_t]
&\le f(x_t)-\eta(1-\beta)\|\nabla f(x_t)\|^2   \\
&\quad +\eta\Bigl[\beta C + (1-\beta)C\Bigr]\|\nabla f(x_t)\|
   +\eta(1-\beta)C\|\nabla f(x_t)\| \\
&\quad +\frac{L}{2}\eta^2 C^2 .
\end{aligned}
\tag{11}
\]
Notice that the two middle lines are identical; summing them yields
\[
\eta C\bigl[\beta + (1-\beta) + (1-\beta)\bigr]\|\nabla f(x_t)\|
= \eta C\bigl[1 + (1-\beta)\bigr]\|\nabla f(x_t)\|
\le 2\eta C\|\nabla f(x_t)\|.
\tag{12}
\]

\noindent\textbf{[Step 8] Use the elementary inequality $2ab\le a^2+b^2$ with $a=\sqrt{1-\beta}\,\|\nabla f(x_t)\|$ and $b=\frac{C}{\sqrt{1-\beta}}$.}
\[
2\eta C\|\nabla f(x_t)\|
\le \eta(1-\beta)\|\nabla f(x_t)\|^2
   +\frac{\eta C^2}{1-\beta}.
\tag{13}
\]

\noindent\textbf{[Step 9] Substitute (13) into (11).}
\[
\mathbb{E}[f(x_{t+1})\mid\mathcal{F}_t]
\le f(x_t)
      -\underbrace{\eta(1-\beta)\|\nabla f(x_t)\|^2}_{\text{descent}}
      +\underbrace{\frac{\eta C^2}{1-\beta}}_{\text{bias term}}
      +\underbrace{\frac{L}{2}\eta^2 C^2}_{\text{smoothness term}} .
\tag{14}
\]

\noindent\textbf{[Step 10] Take full expectation and sum over $t=0,\dots,T-1$.}
\[
\begin{aligned}
\mathbb{E}[f(x_T)] 
&\le f(x_0)
   -\eta(1-\beta)\sum_{t=0}^{T-1}\mathbb{E}\bigl[\|\nabla f(x_t)\|^2\bigr]
   +T\Bigl(\frac{\eta C^2}{1-\beta}
          +\frac{L}{2}\eta^2 C^2\Bigr).
\end{aligned}
\tag{15}
\]

\noindent\textbf{[Step 11] Rearrange (15) and use $f(x_T)\ge f_{\inf}$.}
\[
\eta(1-\beta)\sum_{t=0}^{T-1}\mathbb{E}\bigl[\|\nabla f(x_t)\|^2\bigr]
\le \Delta
   + T\Bigl(\frac{\eta C^2}{1-\beta}
          +\frac{L}{2}\eta^2 C^2\Bigr).
\tag{16}
\]

\noindent\textbf{[Step 12] Divide by $\eta(1-\beta)T$ to obtain the average gradient norm bound.}
\[
\frac{1}{T}\sum_{t=0}^{T-1}\mathbb{E}\bigl[\|\nabla f(x_t)\|^2\bigr]
\le \frac{\Delta}{\eta(1-\beta)T}
   + \frac{C^2}{(1-\beta)^2}
   + \frac{L\eta C^2}{2(1-\beta)} .
\tag{17}
\]

\noindent\textbf{[Step 13] Choose stepsize $\eta\le\frac{1}{L(1+\beta)}$.
Under this restriction the term $\frac{C^2}{(1-\beta)^2}$ can be merged with
$\frac{L\eta C^2}{2(1-\beta)}$ to obtain the cleaner expression stated in the theorem.}
Indeed, using $\eta\le\frac{1}{L(1+\beta)}\le\frac{1}{L}$,
\[
\frac{C^2}{(1-\beta)^2}
\le \frac{L\eta C^2}{2(1-\beta)} + \frac{C^2}{(1-\beta)^2}
\le 2L\eta C^2,
\]
so we can bound the right‑hand side of \eqref{17} by
\[
\frac{2\Delta}{\eta(1-\beta)T}+2L\eta C^2 .
\tag{18}
\]
This is exactly inequality \eqref{1} of Theorem~\ref{thm:main}.

\noindent\textbf{[Step 14] Optimising $\eta$.}
Set $\eta = \sqrt{\frac{2\Delta}{L C^2(1-\beta)T}}$, which respects the stepsize restriction for all sufficiently large $T$.
Plugging into \eqref{18} yields
\[
\frac{1}{T}\sum_{t=0}^{T-1}\mathbb{E}\bigl[\|\nabla f(x_t)\|^2\bigr]
\le 2\sqrt{\frac{2L\Delta C^2}{(1-\beta)T}}
= O\!\bigl(T^{-1/2}\bigr).
\]
Thus the claimed rate follows. \qed

\section*{Rate Derivation (With Constants)}
From inequality \eqref{18} we have for any admissible $\eta$,
\[
\boxed{
\frac{1}{T}\sum_{t=0}^{T-1}\mathbb{E}\bigl[\|\nabla f(x_t)\|^2\bigr]
\le \frac{2\Delta}{\eta(1-\beta)T}+2L\eta C^2 } .
\]
Optimising the RHS over $\eta$ gives the closed‑form optimal stepsize
\[
\eta^\star = \sqrt{\frac{\Delta}{L C^2 (1-\beta) T}} .
\]
Substituting $\eta^\star$ yields the explicit constant:
\[
\frac{1}{T}\sum_{t=0}^{T-1}\mathbb{E}\bigl[\|\nabla f(x_t)\|^2\bigr]
\le 4\sqrt{\frac{L\Delta C^2}{(1-\beta)T}} .
\]

\section*{Special Cases}
\begin{itemize}
\item \textbf{No clipping ($C=\infty$).} Then $\tilde g_t=g_t^{\text{raw}}$ and Lemma~\ref{lem:momentum} no longer guarantees a bound on $\|m_t\|$. The proof collapses because the bias term $2L\eta C^2$ becomes infinite; additional bounded‑gradient assumptions are required.
\item \textbf{Zero momentum ($\beta=0$).} The algorithm reduces to plain clipped SGD. The bound simplifies to
\[
\frac{1}{T}\sum_{t=0}^{T-1}\mathbb{E}\bigl[\|\nabla f(x_t)\|^2\bigr]
\le \frac{2\Delta}{\eta T}+2L\eta C^2,
\]
which is the standard result for SGD with uniformly bounded gradients.
\item \textbf{Deterministic gradients ($\sigma=0$).} The stochastic variance does not appear in the bound; the same $O(T^{-1/2})$ rate holds for the deterministic clipped‑momentum method.
\end{itemize}

\end{document}